\appendix
\section{Baseline Correction: Method Selection}
\label{app:baseline_correction}
\setcounter{table}{0}
\renewcommand{\thetable}{\Alph{section}.\arabic{table}}

We compared four baseline correctors — airPLS \cite{Zhang2010} at the inherited $\lambda=100$, airPLS with $\lambda$ selected per spectrum by generalized cross-validation (GCV) \cite{CravenWahba1979}, airPLS with $\lambda$ selected per spectrum at the L-curve corner \cite{Hansen1992}, and IARPLS \cite{Ye2020} — using a synthetic generator approximating brain-biopsy Raman spectra, plus a ground-truth-free invariance test on real spectra. Four experiments were run in sequence: (1) fitting $\lambda$ for IARPLS, (2) comparing all four correctors against known ground truth, (3) an invariance test on real spectra, (4) IARPLS against an oracle airPLS $\lambda$.

\subsection*{Synthetic Data Generator}
Each generated spectrum is $s(\tilde\nu) = p(\tilde\nu) + b(\tilde\nu) + \varepsilon(\tilde\nu)$ on the real wavenumber axis $\tilde\nu$ (1738 points, $50$--$3400\text{ cm}^{-1}$).

\[
p(\tilde\nu) = \sum_{k=1}^{14} A_k \exp\!\left(-\frac{(\tilde\nu-\mu_k)^2}{2\sigma_k^2}\right), \qquad \sigma_k=\frac{w_k}{2\sqrt{2\ln 2}}
\]
\[
A_k = a\,r_k\,u_k,\quad u_k\sim\mathcal U(0.75,1.25), \qquad w_k = f_k\,v_k,\quad v_k\sim\mathcal U(0.85,1.15), \qquad a=400
\]

\begin{table}[H]
\centering
\small
\caption{Synthetic band positions $\mu_k$, relative amplitudes $r_k$, and FWHM $f_k$.}
\label{tab:baseline_generator_bands}
\begin{tabular*}{\textwidth}{@{\extracolsep{\fill}}lccc}
\toprule
Region & Bands ($\text{cm}^{-1}$) & $r_k$ & $f_k$ ($\text{cm}^{-1}$) \\
\midrule
Protein / nucleic acid & 1004, 1032, 1128, 1230, 1260, 1300, 1340 & $0.30$--$1.00$ & $8$--$18$ \\
Lipid / protein & 1440, 1550, 1660 & $0.35$--$1.20$ & $14$--$20$ \\
Lipid CH stretch & 2850, 2885, 2930 & $1.45$--$1.80$ & $18$--$22$ \\
Aromatic & 3060 & $0.30$ & $20$ \\
\bottomrule
\end{tabular*}
\end{table}

With $t=(\tilde\nu-\tilde\nu_{\min})/\mathrm{ptp}(\tilde\nu)\in[0,1]$,
\[
g(t) = \exp\!\left(-\frac{(t-c)^2}{2w^2}\right) + 0.6\,e^{-t/\tau} + c_1t+c_2t^2+c_3t^3, \qquad b = \alpha\,\frac{g-\min g}{\max(g-\min g)}
\]
with $c\sim\mathcal U(-0.1,0.45)$, $w\sim\mathcal U(0.15,0.45)$, $\tau\sim\mathcal U(0.01,0.04)$, $c_1\sim\mathcal U(-0.3,0.3)$, $c_2\sim\mathcal U(-0.25,0.25)$, $c_3\sim\mathcal U(-0.15,0.15)$, and $\alpha \in \{2,8,25\}\times\mathrm{ptp}(p)$ (low/medium/high fluorescence).
\[
\varepsilon(\tilde\nu)\sim\mathcal N(0,\sigma_n^2), \qquad \sigma_n=\frac{\mathrm{ptp}(p)}{\mathrm{SNR}}, \qquad \mathrm{SNR}=235
\]

\subsection*{Experiment 1: IARPLS $\lambda$ Selection}
Fingerprint-region ($700$--$1800\text{ cm}^{-1}$) NRMSE against $\lambda$, nominal generator, median of 6 spectra per level.

\begin{table}[H]
\centering
\small
\caption{IARPLS NRMSE vs.\ $\lambda$, nominal generator.}
\label{tab:baseline_lambda_fit}
\begin{tabular*}{\textwidth}{@{\extracolsep{\fill}}lccc}
\toprule
$\lambda$ & Low & Medium & High \\
\midrule
$10^2$ & $0.0052$ & $0.0052$ & $0.0052$ \\
$10^3$ & $0.0049$ & $0.0049$ & $0.0049$ \\
$10^4$ & $0.0046$ & $0.0046$ & $0.0046$ \\
$10^5$ & $0.0046$ & $0.0047$ & $0.0269$ \\
$10^6$ & $0.0050$ & $0.1791$ & $2.7035$ \\
$10^7$ & $0.0495$ & $2.7456$ & $3.0016$ \\
\bottomrule
\end{tabular*}
\end{table}

Fingerprint NRMSE vs.\ $\lambda$ across three synthetic baseline curvature families (smooth, structured, sharp), median over lineshape and noise conditions, $n=12$ per cell.

\begin{table}[H]
\centering
\small
\caption{IARPLS NRMSE vs.\ $\lambda$, by baseline curvature family.}
\label{tab:baseline_curvature_family}
\begin{tabular*}{\textwidth}{@{\extracolsep{\fill}}lccc}
\toprule
$\lambda$ & Smooth & Structured & Sharp \\
\midrule
$10^2$ & $0.0169$ & $0.0113$ & $0.0309$ \\
$3.2\times10^2$ & $0.0121$ & $0.0108$ & $0.0290$ \\
$10^3$ & $0.0115$ & $0.0127$ & $0.0400$ \\
$3.2\times10^3$ & $0.0111$ & $0.0114$ & $0.4942$ \\
$10^4$ & $0.0107$ & $0.0122$ & $0.7950$ \\
$10^5$ & $0.0119$ & $0.2313$ & $1.2984$ \\
$10^6$ & $0.8861$ & $1.7073$ & $2.4629$ \\
\bottomrule
\end{tabular*}
\end{table}

\begin{table}[H]
\centering
\small
\caption{Width of the near-optimal $\lambda$ window (within 25\% of the family's minimum NRMSE).}
\label{tab:baseline_curvature_window}
\begin{tabular*}{\textwidth}{@{\extracolsep{\fill}}lcccc}
\toprule
Family & Best $\lambda$ & Best NRMSE & Window & Decades \\
\midrule
Smooth & $10^4$ & $0.0107$ & $3.2\times10^2$--$10^5$ & $2.5$ \\
Structured & $3.2\times10^2$ & $0.0108$ & $10^2$--$3.2\times10^4$ & $2.5$ \\
Sharp & $3.2\times10^2$ & $0.0290$ & $10^2$--$3.2\times10^2$ & $0.5$ \\
\bottomrule
\end{tabular*}
\end{table}

\subsection*{Experiment 2: Comparison Against Ground Truth}
Synthetic spectra, $n=40$ per fluorescence level, $\lambda=10^4$ for IARPLS.

\begin{table}[H]
\centering
\small
\caption{Recovered/true peak height (median over fluorescence levels).}
\label{tab:baseline_peak_fidelity}
\begin{tabular*}{\textwidth}{@{\extracolsep{\fill}}lcccc}
\toprule
Method & 1004 & 1440 & 1660 & 2930 \\
\midrule
airPLS $\lambda=100$ & $0.900$ & $0.691$ & $0.574$ & $0.584$ \\
airPLS+GCV & $0.126$ & $0.025$ & $0.017$ & $0.011$ \\
airPLS+L-curve & $1.002$ & $0.997$ & $0.987$ & $0.927$ \\
IARPLS & $1.003$ & $1.003$ & $1.004$ & $1.003$ \\
\bottomrule
\end{tabular*}
\end{table}

\begin{table}[H]
\centering
\small
\caption{Fingerprint-region NRMSE by fluorescence level.}
\label{tab:baseline_nrmse_level}
\begin{tabular*}{\textwidth}{@{\extracolsep{\fill}}lccc}
\toprule
Method & Low & Medium & High \\
\midrule
airPLS $\lambda=100$ & $0.0267$ & $0.0512$ & $0.0967$ \\
airPLS+GCV & $0.0476$ & $0.1178$ & $0.1134$ \\
airPLS+L-curve & $0.0110$ & $0.0264$ & $0.0245$ \\
IARPLS & $0.0048$ & $0.0047$ & $0.0047$ \\
\bottomrule
\end{tabular*}
\note{GCV pinned $\lambda$ at the grid floor in $83.3\%$ of spectra; L-curve pinned in $0\%$.}
\end{table}

\subsection*{Experiment 3: Invariance on Real Spectra}
Peak-height drift on real spectra when synthetic fluorescence is added ($0$ = invariant); no ground truth required.

\begin{table}[H]
\centering
\small
\caption{Peak-height drift under added fluorescence, real spectra.}
\label{tab:baseline_invariance}
\begin{tabular*}{\textwidth}{@{\extracolsep{\fill}}lccc}
\toprule
Method & Low & Medium & High \\
\midrule
airPLS+GCV & $0.473$ & $0.546$ & $0.848$ \\
airPLS+L-curve & $0.467$ & $0.545$ & $0.632$ \\
airPLS $\lambda=100$ & $0.161$ & $0.274$ & $0.553$ \\
IARPLS & $0.0004$ & $0.0014$ & $0.0030$ \\
\bottomrule
\end{tabular*}
\end{table}

\subsection*{Experiment 4: IARPLS vs.\ Oracle airPLS $\lambda$}
Per-level airPLS $\lambda$ chosen against known ground truth (oracle); fingerprint NRMSE, median of 5 spectra per level.

\begin{table}[H]
\centering
\small
\caption{IARPLS vs.\ oracle airPLS $\lambda$.}
\label{tab:baseline_oracle}
\begin{tabular*}{\textwidth}{@{\extracolsep{\fill}}lccccc}
\toprule
Fluorescence & airPLS $\lambda=100$ & Oracle $\lambda$ & Oracle NRMSE & IARPLS & Ratio \\
\midrule
Low & $0.0219$ & $2.15\times10^4$ & $0.0069$ & $0.0048$ & $1.5\times$ \\
Medium & $0.0430$ & $4.64\times10^4$ & $0.0120$ & $0.0048$ & $2.5\times$ \\
High & $0.0883$ & $2.15\times10^4$ & $0.0141$ & $0.0049$ & $2.9\times$ \\
\bottomrule
\end{tabular*}
\end{table}

IARPLS ($\lambda=10^4$) was adopted as the production baseline corrector. The near-optimal $\lambda$ window narrows from $2.5$ to $0.5$ decades as baseline curvature increases (Table \ref{tab:baseline_curvature_window}); curvature estimated from real spectra exceeds all three synthetic families tested, so $\lambda=10^4$ is retained as a documented default rather than a validated constant. Kizilov et al.\ \cite{Kizilov2026} report a similar $\lambda$-sensitivity for IARPLS on non-biological samples.

\newpage
\section{Extended Experiment Results}
\label{app:extended_results}
\setcounter{table}{0}
\renewcommand{\thetable}{\Alph{section}.\arabic{table}}

This appendix reports the per-experiment, per-model results in full. Each table lists the bag-level accuracy under each training regimen -- one majority-vote decision per bag (biopsy slide), standard metrics over the resulting bag decisions -- the bootstrap difference $\Delta$ relative to the TR/TR baseline with its 95\% confidence interval, and the bag-level Macro F1, Sensitivity, Specificity. Unavailable entries are marked with a dash.

\begin{table}[H]
\centering
\scriptsize
\caption{Extended results for \texttt{exp\_201}: Reconstruction Loss Ablation: MSE-only (gamma=1, delta=0) for binary classification.}
\label{tab:appendix_exp_201}
\begin{tabular*}{\textwidth}{@{\extracolsep{\fill}}llcccccc}
\toprule
\emph{Model} & \emph{Regimen} & \emph{Accuracy} & \emph{$\Delta$} & \emph{95\% CI} & \emph{Macro F1} & \emph{Sensitivity} & \emph{Specificity}\\
\midrule
\multirow{3}{*}{RF + SVM} & TR/TR & $0.741$ & -- & -- & $0.728$ & $0.848$ & $0.600$\\
 & TS/TR & $0.431$ & $-0.310$* & $[-0.414, -0.207]$ & $0.301$ & $0.000$ & $1.000$\\
 & TSR/TR & $0.759$ & $+0.017$ & $[-0.052, +0.086]$ & $0.744$ & $0.879$ & $0.600$\\
\midrule
\multirow{3}{*}{RF} & TR/TR & $0.810$ & -- & -- & $0.797$ & $0.939$ & $0.640$\\
 & TS/TR & $0.707$ & $-0.103$* & $[-0.207, -0.017]$ & $0.696$ & $0.788$ & $0.600$\\
 & TSR/TR & $0.810$ & $+0.000$ & $[+0.000, +0.000]$ & $0.797$ & $0.939$ & $0.640$\\
\midrule
\multirow{3}{*}{XGBoost} & TR/TR & $0.793$ & -- & -- & $0.776$ & $0.939$ & $0.600$\\
 & TS/TR & $0.690$ & $-0.103$ & $[-0.224, +0.017]$ & $0.689$ & $0.576$ & $0.840$\\
 & TSR/TR & $0.810$ & $+0.017$ & $[+0.000, +0.052]$ & $0.797$ & $0.939$ & $0.640$\\
\midrule
\multirow{3}{*}{1D-CNN} & TR/TR & $0.621$ & -- & -- & $0.505$ & $0.970$ & $0.160$\\
 & TS/TR & $0.397$ & $-0.224$* & $[-0.345, -0.103]$ & $0.374$ & $0.182$ & $0.680$\\
 & TSR/TR & $0.603$ & $-0.017$ & $[-0.121, +0.069]$ & $0.556$ & $0.818$ & $0.320$\\
\midrule
CbR & -- & $0.741$ & -- & $[+0.621, +0.845]$ & $0.732$ & $0.818$ & $0.640$\\
\bottomrule
\end{tabular*}
\note{Regimens: TR/TR (real-only baseline), TS/TR (synthetic-only), TSR/TR (synthetic + real), CbR (Classification by Reconstruction, model-independent). Accuracy and all extra columns are bag-level: one majority-vote decision per bag (biopsy slide), standard metrics over the $n$ bag decisions -- see Appendix~\ref{app:baseline_correction}'s sibling methods discussion for the full derivation. $\Delta$ and 95\% CI are the bag-level bootstrap difference vs.\ TR/TR (class-stratified resampling of bag decisions, 1000 iterations); (*) marks a 95\% CI excluding zero. The CbR row reports its own bag-level accuracy 95\% CI (not a difference vs.\ TR/TR), so its $\Delta$ column is left blank.}
\end{table}

\begin{table}[H]
\centering
\scriptsize
\caption{Extended results for \texttt{exp\_202}: Reconstruction Loss Ablation: Balanced MSE+Cosine (gamma=0.5, delta=0.5) for binary classification.}
\label{tab:appendix_exp_202}
\begin{tabular*}{\textwidth}{@{\extracolsep{\fill}}llcccccc}
\toprule
\emph{Model} & \emph{Regimen} & \emph{Accuracy} & \emph{$\Delta$} & \emph{95\% CI} & \emph{Macro F1} & \emph{Sensitivity} & \emph{Specificity}\\
\midrule
\multirow{3}{*}{RF + SVM} & TR/TR & -- & -- & -- & -- & $0.724$ & $0.724$\\
 & TS/TR & -- & -- & -- & -- & $0.500$ & $0.500$\\
 & TSR/TR & -- & -- & -- & -- & $0.739$ & $0.739$\\
\midrule
\multirow{3}{*}{RF} & TR/TR & -- & -- & -- & -- & $0.790$ & $0.790$\\
 & TS/TR & -- & -- & -- & -- & $0.679$ & $0.679$\\
 & TSR/TR & -- & -- & -- & -- & $0.790$ & $0.790$\\
\midrule
\multirow{3}{*}{XGBoost} & TR/TR & -- & -- & -- & -- & $0.770$ & $0.770$\\
 & TS/TR & -- & -- & -- & -- & $0.679$ & $0.679$\\
 & TSR/TR & -- & -- & -- & -- & $0.770$ & $0.770$\\
\midrule
\multirow{3}{*}{1D-CNN} & TR/TR & -- & -- & -- & -- & $0.565$ & $0.565$\\
 & TS/TR & -- & -- & -- & -- & $0.451$ & $0.451$\\
 & TSR/TR & -- & -- & -- & -- & $0.693$ & $0.693$\\
\midrule
CbR & -- & -- & -- & -- & -- & $0.769$ & $0.769$\\
\bottomrule
\end{tabular*}
\note{Regimens: TR/TR (real-only baseline), TS/TR (synthetic-only), TSR/TR (synthetic + real), CbR (Classification by Reconstruction, model-independent). Accuracy and all extra columns are bag-level: one majority-vote decision per bag (biopsy slide), standard metrics over the $n$ bag decisions -- see Appendix~\ref{app:baseline_correction}'s sibling methods discussion for the full derivation. $\Delta$ and 95\% CI are the bag-level bootstrap difference vs.\ TR/TR (class-stratified resampling of bag decisions, 1000 iterations); (*) marks a 95\% CI excluding zero. The CbR row reports its own bag-level accuracy 95\% CI (not a difference vs.\ TR/TR), so its $\Delta$ column is left blank.}
\end{table}

\begin{table}[H]
\centering
\scriptsize
\caption{Extended results for \texttt{exp\_203}: Reconstruction Loss Ablation: Cosine-only (gamma=0, delta=1) for binary classification.}
\label{tab:appendix_exp_203}
\begin{tabular*}{\textwidth}{@{\extracolsep{\fill}}llcccccc}
\toprule
\emph{Model} & \emph{Regimen} & \emph{Accuracy} & \emph{$\Delta$} & \emph{95\% CI} & \emph{Macro F1} & \emph{Sensitivity} & \emph{Specificity}\\
\midrule
\multirow{3}{*}{RF + SVM} & TR/TR & -- & -- & -- & -- & $0.724$ & $0.724$\\
 & TS/TR & -- & -- & -- & -- & $0.500$ & $0.500$\\
 & TSR/TR & -- & -- & -- & -- & $0.719$ & $0.719$\\
\midrule
\multirow{3}{*}{RF} & TR/TR & -- & -- & -- & -- & $0.790$ & $0.790$\\
 & TS/TR & -- & -- & -- & -- & $0.659$ & $0.659$\\
 & TSR/TR & -- & -- & -- & -- & $0.790$ & $0.790$\\
\midrule
\multirow{3}{*}{XGBoost} & TR/TR & -- & -- & -- & -- & $0.770$ & $0.770$\\
 & TS/TR & -- & -- & -- & -- & $0.630$ & $0.630$\\
 & TSR/TR & -- & -- & -- & -- & $0.770$ & $0.770$\\
\midrule
\multirow{3}{*}{1D-CNN} & TR/TR & -- & -- & -- & -- & $0.565$ & $0.565$\\
 & TS/TR & -- & -- & -- & -- & $0.495$ & $0.495$\\
 & TSR/TR & -- & -- & -- & -- & $0.769$ & $0.769$\\
\midrule
CbR & -- & -- & -- & -- & -- & $0.715$ & $0.715$\\
\bottomrule
\end{tabular*}
\note{Regimens: TR/TR (real-only baseline), TS/TR (synthetic-only), TSR/TR (synthetic + real), CbR (Classification by Reconstruction, model-independent). Accuracy and all extra columns are bag-level: one majority-vote decision per bag (biopsy slide), standard metrics over the $n$ bag decisions -- see Appendix~\ref{app:baseline_correction}'s sibling methods discussion for the full derivation. $\Delta$ and 95\% CI are the bag-level bootstrap difference vs.\ TR/TR (class-stratified resampling of bag decisions, 1000 iterations); (*) marks a 95\% CI excluding zero. The CbR row reports its own bag-level accuracy 95\% CI (not a difference vs.\ TR/TR), so its $\Delta$ column is left blank.}
\end{table}

\begin{table}[H]
\centering
\scriptsize
\caption{Extended results for \texttt{exp\_204}: Reconstruction Loss Ablation: MSE-only (gamma=1, delta=0) for 6-class classification.}
\label{tab:appendix_exp_204}
\begin{tabular*}{\textwidth}{@{\extracolsep{\fill}}llcccccc}
\toprule
\emph{Model} & \emph{Regimen} & \emph{Accuracy} & \emph{$\Delta$} & \emph{95\% CI} & \emph{Macro F1} & \emph{Sensitivity} & \emph{Specificity}\\
\midrule
\multirow{3}{*}{RF + SVM} & TR/TR & -- & -- & -- & -- & $0.221$ & $0.851$\\
 & TS/TR & -- & -- & -- & -- & $0.141$ & $0.829$\\
 & TSR/TR & -- & -- & -- & -- & $0.195$ & $0.846$\\
\midrule
\multirow{3}{*}{RF} & TR/TR & -- & -- & -- & -- & $0.224$ & $0.849$\\
 & TS/TR & -- & -- & -- & -- & $0.274$ & $0.853$\\
 & TSR/TR & -- & -- & -- & -- & $0.222$ & $0.849$\\
\midrule
\multirow{3}{*}{XGBoost} & TR/TR & -- & -- & -- & -- & $0.198$ & $0.842$\\
 & TS/TR & -- & -- & -- & -- & $0.289$ & $0.864$\\
 & TSR/TR & -- & -- & -- & -- & $0.198$ & $0.842$\\
\midrule
\multirow{3}{*}{1D-CNN} & TR/TR & -- & -- & -- & -- & $0.151$ & $0.829$\\
 & TS/TR & -- & -- & -- & -- & $0.133$ & $0.829$\\
 & TSR/TR & -- & -- & -- & -- & $0.252$ & $0.851$\\
\midrule
CbR & -- & -- & -- & -- & -- & $0.271$ & $0.852$\\
\bottomrule
\end{tabular*}
\note{Regimens: TR/TR (real-only baseline), TS/TR (synthetic-only), TSR/TR (synthetic + real), CbR (Classification by Reconstruction, model-independent). Accuracy and all extra columns are bag-level: one majority-vote decision per bag (biopsy slide), standard metrics over the $n$ bag decisions -- see Appendix~\ref{app:baseline_correction}'s sibling methods discussion for the full derivation. $\Delta$ and 95\% CI are the bag-level bootstrap difference vs.\ TR/TR (class-stratified resampling of bag decisions, 1000 iterations); (*) marks a 95\% CI excluding zero. The CbR row reports its own bag-level accuracy 95\% CI (not a difference vs.\ TR/TR), so its $\Delta$ column is left blank.}
\end{table}

\begin{table}[H]
\centering
\scriptsize
\caption{Extended results for \texttt{exp\_205}: Reconstruction Loss Ablation: Balanced MSE+Cosine (gamma=0.5, delta=0.5) for 6-class classification.}
\label{tab:appendix_exp_205}
\begin{tabular*}{\textwidth}{@{\extracolsep{\fill}}llcccccc}
\toprule
\emph{Model} & \emph{Regimen} & \emph{Accuracy} & \emph{$\Delta$} & \emph{95\% CI} & \emph{Macro F1} & \emph{Sensitivity} & \emph{Specificity}\\
\midrule
\multirow{3}{*}{RF + SVM} & TR/TR & -- & -- & -- & -- & $0.221$ & $0.851$\\
 & TS/TR & -- & -- & -- & -- & $0.230$ & $0.844$\\
 & TSR/TR & -- & -- & -- & -- & $0.236$ & $0.850$\\
\midrule
\multirow{3}{*}{RF} & TR/TR & -- & -- & -- & -- & $0.227$ & $0.851$\\
 & TS/TR & -- & -- & -- & -- & $0.279$ & $0.852$\\
 & TSR/TR & -- & -- & -- & -- & $0.211$ & $0.846$\\
\midrule
\multirow{3}{*}{XGBoost} & TR/TR & -- & -- & -- & -- & $0.198$ & $0.842$\\
 & TS/TR & -- & -- & -- & -- & $0.281$ & $0.856$\\
 & TSR/TR & -- & -- & -- & -- & $0.183$ & $0.838$\\
\midrule
\multirow{3}{*}{1D-CNN} & TR/TR & -- & -- & -- & -- & $0.174$ & $0.831$\\
 & TS/TR & -- & -- & -- & -- & $0.248$ & $0.843$\\
 & TSR/TR & -- & -- & -- & -- & $0.354$ & $0.872$\\
\midrule
CbR & -- & -- & -- & -- & -- & $0.289$ & $0.854$\\
\bottomrule
\end{tabular*}
\note{Regimens: TR/TR (real-only baseline), TS/TR (synthetic-only), TSR/TR (synthetic + real), CbR (Classification by Reconstruction, model-independent). Accuracy and all extra columns are bag-level: one majority-vote decision per bag (biopsy slide), standard metrics over the $n$ bag decisions -- see Appendix~\ref{app:baseline_correction}'s sibling methods discussion for the full derivation. $\Delta$ and 95\% CI are the bag-level bootstrap difference vs.\ TR/TR (class-stratified resampling of bag decisions, 1000 iterations); (*) marks a 95\% CI excluding zero. The CbR row reports its own bag-level accuracy 95\% CI (not a difference vs.\ TR/TR), so its $\Delta$ column is left blank.}
\end{table}

\begin{table}[H]
\centering
\scriptsize
\caption{Extended results for \texttt{exp\_206}: Reconstruction Loss Ablation: Cosine-only (gamma=0, delta=1) for 6-class classification.}
\label{tab:appendix_exp_206}
\begin{tabular*}{\textwidth}{@{\extracolsep{\fill}}llcccccc}
\toprule
\emph{Model} & \emph{Regimen} & \emph{Accuracy} & \emph{$\Delta$} & \emph{95\% CI} & \emph{Macro F1} & \emph{Sensitivity} & \emph{Specificity}\\
\midrule
\multirow{3}{*}{RF + SVM} & TR/TR & -- & -- & -- & -- & $0.254$ & $0.852$\\
 & TS/TR & -- & -- & -- & -- & $0.161$ & $0.827$\\
 & TSR/TR & -- & -- & -- & -- & $0.258$ & $0.854$\\
\midrule
\multirow{3}{*}{RF} & TR/TR & -- & -- & -- & -- & $0.215$ & $0.844$\\
 & TS/TR & -- & -- & -- & -- & $0.213$ & $0.842$\\
 & TSR/TR & -- & -- & -- & -- & $0.208$ & $0.843$\\
\midrule
\multirow{3}{*}{XGBoost} & TR/TR & -- & -- & -- & -- & $0.215$ & $0.845$\\
 & TS/TR & -- & -- & -- & -- & $0.207$ & $0.843$\\
 & TSR/TR & -- & -- & -- & -- & $0.213$ & $0.844$\\
\midrule
\multirow{3}{*}{1D-CNN} & TR/TR & -- & -- & -- & -- & $0.168$ & $0.834$\\
 & TS/TR & -- & -- & -- & -- & $0.109$ & $0.825$\\
 & TSR/TR & -- & -- & -- & -- & $0.263$ & $0.856$\\
\midrule
CbR & -- & -- & -- & -- & -- & $0.197$ & $0.842$\\
\bottomrule
\end{tabular*}
\note{Regimens: TR/TR (real-only baseline), TS/TR (synthetic-only), TSR/TR (synthetic + real), CbR (Classification by Reconstruction, model-independent). Accuracy and all extra columns are bag-level: one majority-vote decision per bag (biopsy slide), standard metrics over the $n$ bag decisions -- see Appendix~\ref{app:baseline_correction}'s sibling methods discussion for the full derivation. $\Delta$ and 95\% CI are the bag-level bootstrap difference vs.\ TR/TR (class-stratified resampling of bag decisions, 1000 iterations); (*) marks a 95\% CI excluding zero. The CbR row reports its own bag-level accuracy 95\% CI (not a difference vs.\ TR/TR), so its $\Delta$ column is left blank.}
\end{table}

\begin{table}[H]
\centering
\scriptsize
\caption{Extended results for \texttt{exp\_302}: Latent Space Dimension Ablation: d=16 (MSE-only, binary classification).}
\label{tab:appendix_exp_302}
\begin{tabular*}{\textwidth}{@{\extracolsep{\fill}}llcccccc}
\toprule
\emph{Model} & \emph{Regimen} & \emph{Accuracy} & \emph{$\Delta$} & \emph{95\% CI} & \emph{Macro F1} & \emph{Sensitivity} & \emph{Specificity}\\
\midrule
\multirow{3}{*}{RF + SVM} & TR/TR & -- & -- & -- & -- & $0.724$ & $0.724$\\
 & TS/TR & -- & -- & -- & -- & $0.511$ & $0.511$\\
 & TSR/TR & -- & -- & -- & -- & $0.739$ & $0.739$\\
\midrule
\multirow{3}{*}{RF} & TR/TR & -- & -- & -- & -- & $0.790$ & $0.790$\\
 & TS/TR & -- & -- & -- & -- & $0.699$ & $0.699$\\
 & TSR/TR & -- & -- & -- & -- & $0.775$ & $0.775$\\
\midrule
\multirow{3}{*}{XGBoost} & TR/TR & -- & -- & -- & -- & $0.770$ & $0.770$\\
 & TS/TR & -- & -- & -- & -- & $0.735$ & $0.735$\\
 & TSR/TR & -- & -- & -- & -- & $0.790$ & $0.790$\\
\midrule
\multirow{3}{*}{1D-CNN} & TR/TR & -- & -- & -- & -- & $0.565$ & $0.565$\\
 & TS/TR & -- & -- & -- & -- & $0.431$ & $0.431$\\
 & TSR/TR & -- & -- & -- & -- & $0.615$ & $0.615$\\
\midrule
CbR & -- & -- & -- & -- & -- & $0.709$ & $0.709$\\
\bottomrule
\end{tabular*}
\note{Regimens: TR/TR (real-only baseline), TS/TR (synthetic-only), TSR/TR (synthetic + real), CbR (Classification by Reconstruction, model-independent). Accuracy and all extra columns are bag-level: one majority-vote decision per bag (biopsy slide), standard metrics over the $n$ bag decisions -- see Appendix~\ref{app:baseline_correction}'s sibling methods discussion for the full derivation. $\Delta$ and 95\% CI are the bag-level bootstrap difference vs.\ TR/TR (class-stratified resampling of bag decisions, 1000 iterations); (*) marks a 95\% CI excluding zero. The CbR row reports its own bag-level accuracy 95\% CI (not a difference vs.\ TR/TR), so its $\Delta$ column is left blank.}
\end{table}

\begin{table}[H]
\centering
\scriptsize
\caption{Extended results for \texttt{exp\_303}: Latent Space Dimension Ablation: d=64 (MSE-only, binary classification).}
\label{tab:appendix_exp_303}
\begin{tabular*}{\textwidth}{@{\extracolsep{\fill}}llcccccc}
\toprule
\emph{Model} & \emph{Regimen} & \emph{Accuracy} & \emph{$\Delta$} & \emph{95\% CI} & \emph{Macro F1} & \emph{Sensitivity} & \emph{Specificity}\\
\midrule
\multirow{3}{*}{RF + SVM} & TR/TR & -- & -- & -- & -- & $0.724$ & $0.724$\\
 & TS/TR & -- & -- & -- & -- & $0.500$ & $0.500$\\
 & TSR/TR & -- & -- & -- & -- & $0.759$ & $0.759$\\
\midrule
\multirow{3}{*}{RF} & TR/TR & -- & -- & -- & -- & $0.775$ & $0.775$\\
 & TS/TR & -- & -- & -- & -- & $0.659$ & $0.659$\\
 & TSR/TR & -- & -- & -- & -- & $0.775$ & $0.775$\\
\midrule
\multirow{3}{*}{XGBoost} & TR/TR & -- & -- & -- & -- & $0.770$ & $0.770$\\
 & TS/TR & -- & -- & -- & -- & $0.593$ & $0.593$\\
 & TSR/TR & -- & -- & -- & -- & $0.790$ & $0.790$\\
\midrule
\multirow{3}{*}{1D-CNN} & TR/TR & -- & -- & -- & -- & $0.585$ & $0.585$\\
 & TS/TR & -- & -- & -- & -- & $0.500$ & $0.500$\\
 & TSR/TR & -- & -- & -- & -- & $0.735$ & $0.735$\\
\midrule
CbR & -- & -- & -- & -- & -- & $0.664$ & $0.664$\\
\bottomrule
\end{tabular*}
\note{Regimens: TR/TR (real-only baseline), TS/TR (synthetic-only), TSR/TR (synthetic + real), CbR (Classification by Reconstruction, model-independent). Accuracy and all extra columns are bag-level: one majority-vote decision per bag (biopsy slide), standard metrics over the $n$ bag decisions -- see Appendix~\ref{app:baseline_correction}'s sibling methods discussion for the full derivation. $\Delta$ and 95\% CI are the bag-level bootstrap difference vs.\ TR/TR (class-stratified resampling of bag decisions, 1000 iterations); (*) marks a 95\% CI excluding zero. The CbR row reports its own bag-level accuracy 95\% CI (not a difference vs.\ TR/TR), so its $\Delta$ column is left blank.}
\end{table}

\begin{table}[H]
\centering
\scriptsize
\caption{Extended results for \texttt{exp\_401}: TSRTR ratio ablation 1:0.5 — real+synth enriched (reuses exp\_201 CVAE/TRTR).}
\label{tab:appendix_exp_401}
\begin{tabular*}{\textwidth}{@{\extracolsep{\fill}}llcccccc}
\toprule
\emph{Model} & \emph{Regimen} & \emph{Accuracy} & \emph{$\Delta$} & \emph{95\% CI} & \emph{Macro F1} & \emph{Sensitivity} & \emph{Specificity}\\
\midrule
\multirow{3}{*}{RF + SVM} & TR/TR & $0.741$ & -- & -- & $0.728$ & $0.848$ & $0.600$\\
 & TS/TR & -- & -- & -- & -- & $0.500$ & $0.500$\\
 & TSR/TR & $0.724$ & $-0.017$ & $[-0.082, +0.048]$ & $0.712$ & $0.818$ & $0.600$\\
\midrule
\multirow{3}{*}{RF} & TR/TR & $0.810$ & -- & -- & $0.797$ & $0.939$ & $0.640$\\
 & TS/TR & -- & -- & -- & -- & $0.694$ & $0.694$\\
 & TSR/TR & $0.793$ & $-0.017$ & $[-0.065, +0.000]$ & $0.780$ & $0.909$ & $0.640$\\
\midrule
\multirow{3}{*}{XGBoost} & TR/TR & $0.793$ & -- & -- & $0.776$ & $0.939$ & $0.600$\\
 & TS/TR & -- & -- & -- & -- & $0.708$ & $0.708$\\
 & TSR/TR & $0.793$ & $+0.000$ & $[+0.000, +0.000]$ & $0.776$ & $0.939$ & $0.600$\\
\midrule
\multirow{3}{*}{1D-CNN} & TR/TR & $0.621$ & -- & -- & $0.505$ & $0.970$ & $0.160$\\
 & TS/TR & -- & -- & -- & -- & $0.431$ & $0.431$\\
 & TSR/TR & $0.672$ & $+0.052$ & $[-0.048, +0.182]$ & $0.670$ & $0.667$ & $0.680$\\
\midrule
CbR & -- & $0.741$ & -- & $[+0.655, +0.845]$ & $0.732$ & $0.818$ & $0.640$\\
\bottomrule
\end{tabular*}
\note{Regimens: TR/TR (real-only baseline), TS/TR (synthetic-only), TSR/TR (synthetic + real), CbR (Classification by Reconstruction, model-independent). Accuracy and all extra columns are bag-level: one majority-vote decision per bag (biopsy slide), standard metrics over the $n$ bag decisions -- see Appendix~\ref{app:baseline_correction}'s sibling methods discussion for the full derivation. $\Delta$ and 95\% CI are the bag-level bootstrap difference vs.\ TR/TR (class-stratified resampling of bag decisions, 1000 iterations); (*) marks a 95\% CI excluding zero. The CbR row reports its own bag-level accuracy 95\% CI (not a difference vs.\ TR/TR), so its $\Delta$ column is left blank.}
\end{table}

\begin{table}[H]
\centering
\scriptsize
\caption{Extended results for \texttt{exp\_402}: TSRTR ratio ablation 1:1 — real+synth enriched (reuses exp\_201 CVAE/TRTR).}
\label{tab:appendix_exp_402}
\begin{tabular*}{\textwidth}{@{\extracolsep{\fill}}llcccccc}
\toprule
\emph{Model} & \emph{Regimen} & \emph{Accuracy} & \emph{$\Delta$} & \emph{95\% CI} & \emph{Macro F1} & \emph{Sensitivity} & \emph{Specificity}\\
\midrule
\multirow{3}{*}{RF + SVM} & TR/TR & -- & -- & -- & -- & $0.724$ & $0.724$\\
 & TS/TR & -- & -- & -- & -- & $0.500$ & $0.500$\\
 & TSR/TR & -- & -- & -- & -- & $0.739$ & $0.739$\\
\midrule
\multirow{3}{*}{RF} & TR/TR & -- & -- & -- & -- & $0.790$ & $0.790$\\
 & TS/TR & -- & -- & -- & -- & $0.694$ & $0.694$\\
 & TSR/TR & -- & -- & -- & -- & $0.775$ & $0.775$\\
\midrule
\multirow{3}{*}{XGBoost} & TR/TR & -- & -- & -- & -- & $0.770$ & $0.770$\\
 & TS/TR & -- & -- & -- & -- & $0.708$ & $0.708$\\
 & TSR/TR & -- & -- & -- & -- & $0.790$ & $0.790$\\
\midrule
\multirow{3}{*}{1D-CNN} & TR/TR & -- & -- & -- & -- & $0.565$ & $0.565$\\
 & TS/TR & -- & -- & -- & -- & $0.431$ & $0.431$\\
 & TSR/TR & -- & -- & -- & -- & $0.673$ & $0.673$\\
\midrule
CbR & -- & -- & -- & -- & -- & $0.729$ & $0.729$\\
\bottomrule
\end{tabular*}
\note{Regimens: TR/TR (real-only baseline), TS/TR (synthetic-only), TSR/TR (synthetic + real), CbR (Classification by Reconstruction, model-independent). Accuracy and all extra columns are bag-level: one majority-vote decision per bag (biopsy slide), standard metrics over the $n$ bag decisions -- see Appendix~\ref{app:baseline_correction}'s sibling methods discussion for the full derivation. $\Delta$ and 95\% CI are the bag-level bootstrap difference vs.\ TR/TR (class-stratified resampling of bag decisions, 1000 iterations); (*) marks a 95\% CI excluding zero. The CbR row reports its own bag-level accuracy 95\% CI (not a difference vs.\ TR/TR), so its $\Delta$ column is left blank.}
\end{table}

\begin{table}[H]
\centering
\scriptsize
\caption{Extended results for \texttt{exp\_403}: TSRTR ratio ablation 1:2 — real+synth enriched (reuses exp\_201 CVAE/TRTR).}
\label{tab:appendix_exp_403}
\begin{tabular*}{\textwidth}{@{\extracolsep{\fill}}llcccccc}
\toprule
\emph{Model} & \emph{Regimen} & \emph{Accuracy} & \emph{$\Delta$} & \emph{95\% CI} & \emph{Macro F1} & \emph{Sensitivity} & \emph{Specificity}\\
\midrule
\multirow{3}{*}{RF + SVM} & TR/TR & -- & -- & -- & -- & $0.724$ & $0.724$\\
 & TS/TR & -- & -- & -- & -- & $0.500$ & $0.500$\\
 & TSR/TR & -- & -- & -- & -- & $0.759$ & $0.759$\\
\midrule
\multirow{3}{*}{RF} & TR/TR & -- & -- & -- & -- & $0.790$ & $0.790$\\
 & TS/TR & -- & -- & -- & -- & $0.694$ & $0.694$\\
 & TSR/TR & -- & -- & -- & -- & $0.790$ & $0.790$\\
\midrule
\multirow{3}{*}{XGBoost} & TR/TR & -- & -- & -- & -- & $0.770$ & $0.770$\\
 & TS/TR & -- & -- & -- & -- & $0.708$ & $0.708$\\
 & TSR/TR & -- & -- & -- & -- & $0.790$ & $0.790$\\
\midrule
\multirow{3}{*}{1D-CNN} & TR/TR & -- & -- & -- & -- & $0.565$ & $0.565$\\
 & TS/TR & -- & -- & -- & -- & $0.431$ & $0.431$\\
 & TSR/TR & -- & -- & -- & -- & $0.755$ & $0.755$\\
\midrule
CbR & -- & -- & -- & -- & -- & $0.729$ & $0.729$\\
\bottomrule
\end{tabular*}
\note{Regimens: TR/TR (real-only baseline), TS/TR (synthetic-only), TSR/TR (synthetic + real), CbR (Classification by Reconstruction, model-independent). Accuracy and all extra columns are bag-level: one majority-vote decision per bag (biopsy slide), standard metrics over the $n$ bag decisions -- see Appendix~\ref{app:baseline_correction}'s sibling methods discussion for the full derivation. $\Delta$ and 95\% CI are the bag-level bootstrap difference vs.\ TR/TR (class-stratified resampling of bag decisions, 1000 iterations); (*) marks a 95\% CI excluding zero. The CbR row reports its own bag-level accuracy 95\% CI (not a difference vs.\ TR/TR), so its $\Delta$ column is left blank.}
\end{table}


\section{Stratified Evaluation Protocol}
\label{app:cv_schema}
\setcounter{table}{0}
\renewcommand{\thetable}{\Alph{section}.\arabic{table}}

We describe here in detail the five-fold evaluation protocol that we replicated from Lita et al. \cite{apollonov2024raman}. Due to the imbalanced nature of the dataset (See Table \ref{tab:dataset}), a strict CV partition stratified at the patient level would result in folds where classes LGm1 or LGm6 would be absent from the test set of one fold. Therefore, assignments were determined manually to ensure that each class contributes at least 15\% of its samples to each fold's test partition; details are listed below. As a consequence, slides from patients in underrepresented classes appear in the test partitions of two folds and in the training partitions of the remaining three; all other slides appear in exactly one test partition. The complete assignment is stored in \texttt{data/splits/manual\_splits.json}; training sets are not the complement of the listed test sets and cannot be reconstructed from the test listings alone. Refer to the \texttt{manual\_splits.json} file for details.
  All tissue sections acquired from the same patient are kept together and assigned exclusively to either the training or the test partition within a given fold, preventing patient-level data leakage.

\paragraph{File structure.} Each key \texttt{"fold\_\textit{k}"} ($k = 0, \ldots, 4$) maps to an object containing two string arrays:

\begin{verbatim}
{
  "fold_0": {
    "train": ["HF-<patient>_<label>.h5_<cluster>", ...],
    "test":  ["HF-<patient>_<label>.h5_<cluster>", ...]
  },
  ...
  "fold_4": {
    "train": [...],
    "test":  [...]
  }
}
\end{verbatim}

Each slide identifier takes the form \texttt{HF-\textit{N}\_\textit{label}.h5}, where \texttt{N} is the patient case number and \texttt{label} identifies the tissue block and acquisition section.

\begin{table}[H]
\centering
\small
\caption{Stratified fold sizes. Training sets are not the complement of the test sets; see text.}
\label{tab:cv_folds_summary}
\begin{tabular}{lrr}
\toprule
Fold & Train slides & Test slides \\
\midrule
0 & 44 & 14 \\
1 & 46 & 12 \\
2 & 44 & 14 \\
3 & 48 & 10 \\
4 & 46 & 12 \\
\bottomrule
\end{tabular}
\note{Counts refer to individual slide acquisitions, not patients. Four slides appear in the test partitions of both fold~0 and fold~4 (patients \texttt{HF-2104} and \texttt{HF-592}), so the sum of test-set sizes across all five folds (62) exceeds the number of unique slides in the dataset (58).}
\end{table}

\paragraph{Slides appearing in two test folds.}
Four slides appear in the test partitions of both fold~0 and fold~4 and are therefore included in the training partitions of folds~1, 2, and~3 only: \texttt{HF-2104\_\#5\_1.h5\_0}, \texttt{HF-2104\_\#9\_1.h5\_1}, and \texttt{HF-2104\_V1T\_1.h5\_2} (patient~\texttt{HF-2104}), and \texttt{HF-592\_V3T\_1.h5\_4} (patient~\texttt{HF-592}). Their repeated appearance can be verified in the listings below.

\paragraph{Test-set assignments.} The slide identifiers held out in each fold are listed verbatim below. Full training-set listings are in the referenced JSON file.

\medskip
\noindent\textbf{Fold 0} (14 slides):
\begin{verbatim}
HF-448_V5B_1.h5_3        HF-305_v4b_1_1.h5_6      HF-615_V5BB_1.h5_9
HF-2104_#5_1.h5_0        HF-2104_#9_1.h5_1         HF-2104_V1T_1.h5_2
HF-442_V4BB_1.h5_12      HF-1002_V1AT_1.h5_0       HF-1002_V2AT_1.h5_1
HF-682_V3AT_1.h5_9       HF-682_V3BB_1.h5_10       HF-894_9_1.h5_11
HF-894_V1BB_1.h5_12      HF-592_V3T_1.h5_4
\end{verbatim}

\medskip
\noindent\textbf{Fold 1} (12 slides):
\begin{verbatim}
HF-868_1_2.h5_4          HF-901_V2T_2.h5_10        HF-960_VIAT_2.h5_11
HF-2614_V1B_1.h5_3       HF-1825_V2B_1.h5_2        HF-2102_V2BB_1.h5_3
HF-2102_V3AM_1.h5_4      HF-2102_V3AM_2.h5_5       HF-988_V1-T_1.h5_13
HF-988_V1B_1.h5_14       HF-1043_V1AM_1.h5_0       HF-2106_V3AM_1.h5_0
\end{verbatim}

\medskip
\noindent\textbf{Fold 2} (14 slides):
\begin{verbatim}
HF-1293_13_1.h5_0        HF-1010_V1T_1.h5_0        HF-1016_IAT_2.h5_1
HF-1334_V58-B_2_1.h5_2   HF-2849_VIT2_1.h5_4       HF-2849_VIT2_1.h5_5
HF-2849_VIT2_2.h5_6      HF-2849_VIT_2_new2021.h5_7 HF-2454_V1AT_1.h5_6
HF-2548_V1T_1.h5_7       HF-1086_#1_1.h5_1         HF-2355_V2AM_1.h5_2
HF-2493_V1T_1.h5_1       HF-2493_V1T_2.h5_2
\end{verbatim}

\medskip
\noindent\textbf{Fold 3} (10 slides):
\begin{verbatim}
HF-1295_V3AM_2.h5_1      HF-2070_V1T_1.h5_4        HF-2776_V2B_2.h5_5
HF-2852_VIT_2_2.h5_8     HF-2715_VIL_1.h5_8        HF-2802_V3T_1.h5_9
HF-2485_V1B_1.h5_3       HF-2600_V1B_1.h5_4        HF-2608_V1T_1.h5_5
HF-2544_V1B_1.h5_3
\end{verbatim}

\medskip
\noindent\textbf{Fold 4} (12 slides):
\begin{verbatim}
HF-2534_V2B_1.h5_2       HF-3271_VIB_2.h5_7        HF-3337_V3T_1.h5_8
HF-2104_#5_1.h5_0        HF-2104_#9_1.h5_1         HF-2104_V1T_1.h5_2
HF-2876_V1T_1.h5_10      HF-2898_V1T_1.h5_11       HF-2619_V1T_1.h5_6
HF-2619_V4T_1.h5_7       HF-2666_V2B_1.h5_8        HF-592_V3T_1.h5_4
\end{verbatim}
